% !TEX root = ../main.tex
\section{Experimental Setup}
We train NS-SDN under a one-step-ahead conditional forecasting regime:
\begin{itemize}
\item Daily DGS10 observations (log-transformed),
\item Single train/holdout split, with the final 365 observations reserved as a holdout window,
\item Teacher forcing during training and conditional one-step evaluation \cite{flunkertdeepar},
\item $K=16$ spectral components (four representative components shown in figures),
\item Latent state dimension $d = 128$, GRU transition.
\end{itemize}
Training uses Adam \cite{adamoptim} (two-stage: $\mathrm{lr}=10^{-3}$ for 1{,}200 steps, then $\mathrm{lr}=4\!\times\!10^{-4}$ for 1{,}800 steps) with smoothness penalties on amplitude and frequency drift. The focus is not forecasting accuracy alone, but the behavior of:
\begin{itemize}
\item Instantaneous frequencies $\omega_{t,k}$,
\item Amplitudes $A_{t,k}$,
\item Gating weights $g_{t,k}$.
\end{itemize}
